% Options for packages loaded elsewhere
\PassOptionsToPackage{unicode}{hyperref}
\PassOptionsToPackage{hyphens}{url}
\documentclass[
  ignorenonframetext,
]{beamer}
\newif\ifbibliography
\usepackage{pgfpages}
\setbeamertemplate{caption}[numbered]
\setbeamertemplate{caption label separator}{: }
\setbeamercolor{caption name}{fg=normal text.fg}
\beamertemplatenavigationsymbolsempty
% remove section numbering
\setbeamertemplate{part page}{
  \centering
  \begin{beamercolorbox}[sep=16pt,center]{part title}
    \usebeamerfont{part title}\insertpart\par
  \end{beamercolorbox}
}
\setbeamertemplate{section page}{
  \centering
  \begin{beamercolorbox}[sep=12pt,center]{section title}
    \usebeamerfont{section title}\insertsection\par
  \end{beamercolorbox}
}
\setbeamertemplate{subsection page}{
  \centering
  \begin{beamercolorbox}[sep=8pt,center]{subsection title}
    \usebeamerfont{subsection title}\insertsubsection\par
  \end{beamercolorbox}
}
% Prevent slide breaks in the middle of a paragraph
\widowpenalties 1 10000
\raggedbottom
\AtBeginPart{
  \frame{\partpage}
}
\AtBeginSection{
  \ifbibliography
  \else
    \frame{\sectionpage}
  \fi
}
\AtBeginSubsection{
  \frame{\subsectionpage}
}
\usepackage{iftex}
\ifPDFTeX
  \usepackage[T1]{fontenc}
  \usepackage[utf8]{inputenc}
  \usepackage{textcomp} % provide euro and other symbols
\else % if luatex or xetex
  \usepackage{unicode-math} % this also loads fontspec
  \defaultfontfeatures{Scale=MatchLowercase}
  \defaultfontfeatures[\rmfamily]{Ligatures=TeX,Scale=1}
\fi
\usepackage{lmodern}
\ifPDFTeX\else
  % xetex/luatex font selection
\fi
% Use upquote if available, for straight quotes in verbatim environments
\IfFileExists{upquote.sty}{\usepackage{upquote}}{}
\IfFileExists{microtype.sty}{% use microtype if available
  \usepackage[]{microtype}
  \UseMicrotypeSet[protrusion]{basicmath} % disable protrusion for tt fonts
}{}
\makeatletter
\@ifundefined{KOMAClassName}{% if non-KOMA class
  \IfFileExists{parskip.sty}{%
    \usepackage{parskip}
  }{% else
    \setlength{\parindent}{0pt}
    \setlength{\parskip}{6pt plus 2pt minus 1pt}}
}{% if KOMA class
  \KOMAoptions{parskip=half}}
\makeatother
\setlength{\emergencystretch}{3em} % prevent overfull lines
\providecommand{\tightlist}{%
  \setlength{\itemsep}{0pt}\setlength{\parskip}{0pt}}
\usepackage{bookmark}
\IfFileExists{xurl.sty}{\usepackage{xurl}}{} % add URL line breaks if available
\urlstyle{same}
\hypersetup{
  hidelinks,
  pdfcreator={LaTeX via pandoc}}

\author{\texorpdfstring{}{}}
\date{}

\begin{document}

\section{\texorpdfstring{\emph{America A Prophecy} and the Political
Economy of Revolutionary
Perception}{America A Prophecy and the Political Economy of Revolutionary Perception}}\label{america-a-prophecy-and-the-political-economy-of-revolutionary-perception}

\begin{frame}{``Shaking their mental chains'': Blake's Prophecy at the
Monetary Threshold}
\protect\phantomsection\label{shaking-their-mental-chains-blakes-prophecy-at-the-monetary-threshold}
\emph{America A Prophecy} (1793) is structured in two movements: a
Preludium of mythological unchaining and a Prophecy of revolutionary
confrontation. At its center lies a passage of economic and
epistemological resonance. When the Thirteen Governors convene in
Bernard's house and the flames of Orc engulf the land, the colonial
administrators react with terror:

\begin{quote}
``Shaking their mental chains, they rush in fury to the sea To quench
their anguish; at the feet of Washington down fall'n They grovel on the
sand and writhing lie'' (\emph{America A Prophecy}, Plates 13--14;
Erdman p.~55) {[}@erdman1988complete; @blake1793america{]}
\end{quote}

The phrase ``mental chains'' is Blake's articulation of the relationship
between political subjugation and epistemic imprisonment. These are not
physical manacles but conceptual constraints---the same ``mind-forg'd
manacles'' that Blake would later identify in \emph{London} (Songs of
Experience, 1794; Erdman p.~70) {[}@erdman1988complete{]}. In the
context of the bimetallic monetary crisis, the ``mental chains'' are the
very statutory ratios, fixed prices, and abstract ledgers that bind
human labor to metallic equivalences. The governors do not merely lose
political authority; they lose the entire conceptual apparatus through
which they have administered value.
\end{frame}

\begin{frame}{Boston's Angel and the Commodification of Human Virtue}
\protect\phantomsection\label{bostons-angel-and-the-commodification-of-human-virtue}
The most economically resonant passage in \emph{America} is Boston's
Angel's speech, which Blake positions as the poem's intellectual climax:

\begin{quote}
``Who commanded this? What God? What Angel? To keep the gen'rous from
experience till the ungenerous Are unrestrain'd performers of the
energies of nature; Till pity is become a trade, and generosity a
science That men get rich by; and the sandy desert is giv'n to the
strong?'' (\emph{America A Prophecy}, Plate 9; Erdman p.~53)
{[}@erdman1988complete; @blake1793america{]}
\end{quote}

This passage demands economic exegesis. ``Pity is become a trade''
anticipates by eighty years Karl Marx's analysis of how capitalism
transforms all human relations into commodity relations: even compassion
itself becomes a tradeable asset, a resource to be deployed for profit.
``Generosity a science / That men get rich by'' describes with precision
the logic of speculative philanthropy, charitable endowments, and the
broader financialization of human virtue that was already visible in the
operations of the Bank of England's credit system
{[}@poovey2008genres{]}. Patrick Brantlinger's \emph{Fictions of State}
traces how this transformation of moral sentiment into financial
infrastructure extended through the entire edifice of 18th-century
credit culture, creating what he terms a ``culture of speculation'' in
which even patriotism and public spirit were securitized
{[}@brantlinger1996fictions{]}.

As Saree Makdisi argues, Blake's \emph{America} challenges the emerging
commodity-form of perception itself: the idea that knowledge, like
goods, is something passively received from external authority rather
than actively produced through imaginative engagement
{[}@makdisi2003william{]}. Boston's Angel demands to know \emph{who has
the right to restrict experience}---a question that applies with equal
force to the sovereign's right to fix the price of metal and the Bank's
right to restrict the conversion of paper into gold. Jon Mee's
\emph{Dangerous Enthusiasm} contextualizes Blake within the broader
culture of 1790s radical dissent, demonstrating that Blake's critique of
commodified experience drew directly on the millenarian and antinomian
traditions that informed popular resistance to enclosure, taxation, and
the disciplinary apparatus of Pitt's wartime state
{[}@mee1992dangerous{]}.
\end{frame}

\begin{frame}{The ``Crime'' of 1873 and the Demonetization of Silver}
\protect\phantomsection\label{the-crime-of-1873-and-the-demonetization-of-silver}
Blake's prophetic framework proves its relevance when extended beyond
its immediate historical context to encompass the full arc of American
bimetallic politics. The Fourth Coinage Act of \textbf{12 February
1873}---known to its critics as ``The Crime of 1873''---demonetized the
silver dollar by omitting it from the list of authorized coins,
effectively placing the United States on a gold monometallic standard
{[}@laughlin1898bimetallism{]}. This legislative act, passed with little
public debate while the nation was still recovering from Civil War and
the Reconstruction crisis, enacted on American soil the same regression
that Blake had diagnosed in Britain's 1816 Coinage Act. As Richard White
documents in \emph{The Republic for Which It Stands}, the 1873 Act was
not a sudden conspiracy but the culmination of a decade-long lobbying
campaign by Eastern financial interests who sought to align American
monetary policy with the emerging international gold consensus centered
on London and Berlin {[}@white2017republic{]}.

The parallels with the British experience are structurally exact. Where
Lord Liverpool's 1816 Act (56 Geo. III c.68) demonetized full-bodied
silver by reducing it to token coinage with a 40-shilling legal tender
limit, the 1873 Act accomplished the same amputation by erasing the
silver dollar from the statute book. As Marc Flandreau argues in
\emph{The Glitter of Gold}, this global shift from bimetallism to gold
was fundamentally a struggle over \emph{political theology}---a contest
between competing visions of what money \emph{is} and whom it should
serve {[}@flandreau2004glitter{]}. Agrarian producers and Western miners
defended silver as a tangible sacrament of their labor, while Eastern
financial elites demanded the hierarchical, centralized abstraction of
the gold standard. Richard Franklin Bensel's \emph{Political Economy of
American Industrialization} demonstrates that the deflation following
1873 was not an incidental consequence but a structural requirement of
the gold standard regime, which transferred wealth systematically from
debtor agrarians to creditor financiers {[}@bensel2000political{]}.

Where Newton's Mint ratio had overvalued gold to drain England of
silver, the post-1873 deflation---in which Milton Friedman notes
American price levels collapsed by roughly 50\% between 1865 and
1896---crushed American farmers under a contracting gold supply
{[}@friedman1963monetary{]}. The producing classes of both nations---the
``laborious'' upon whom Blake's ``Councellor'' throws the ``curb / Of
Poverty'' (\emph{Jerusalem}, Plate 27; Erdman p.~170)---bore the
consequences of this monetary contraction, while the urban creditor
class reaped profits from the rising value of their gold-denominated
claims {[}@delmar1895history; @erdman1988complete{]}.
\end{frame}

\begin{frame}{The Free Silver Movement and the Populist Insurgency}
\protect\phantomsection\label{the-free-silver-movement-and-the-populist-insurgency}
The political response to the Crime of 1873 constituted nothing less
than the most significant challenge to financial orthodoxy in American
history. The Free Silver movement (1873--1900) demanded the restoration
of free coinage of silver at a ratio of \textbf{16:1}---a deliberate
rejection of the monometallic stranglehold that had impoverished the
agrarian West and South {[}@laughlin1898bimetallism{]}. Charles Postel's
\emph{The Populist Vision} demonstrates that the Populists were not, as
Richard Hofstadter once dismissively characterized them,
backward-looking agrarian romantics, but sophisticated political actors
who understood the monetary system's structural biases with considerable
analytical precision {[}@postel2007populist; @hofstadter1955age{]}.
Congress responded with the \textbf{Bland-Allison Act of 1878},
requiring the Treasury to purchase \$2--4 million of silver bullion
monthly for coinage, and the \textbf{Sherman Silver Purchase Act of
1890}, mandating the purchase of 4.5 million ounces of silver per
month---legislative interventions that represented partial attempts to
restore the bimetallic synthesis against the crushing monometallism
imposed by Eastern financial interests.

The movement reached its crescendo in William Jennings Bryan's
\textbf{``Cross of Gold'' speech} at the 1896 Democratic National
Convention in Chicago (9 July 1896). As Michael Kazin documents in
\emph{A Godly Hero}, Bryan's address was not merely a political
performance but a distillation of two decades of agrarian rage against
the deflationary regime imposed by the gold standard
{[}@kazin2006godly{]}. His defense of the agrarian producing classes
against urban financial elites---``You come to us and tell us that the
great cities are in favor of the gold standard; we reply that the great
cities rest upon our broad and fertile prairies. Burn down your cities
and leave our farms, and your cities will spring up again as if by
magic; but destroy our farms and the grass will grow in the streets of
every city in the country''---recapitulates with precision Blake's
insistence that productive, embodied labor constitutes the only true
source of value {[}@laughlin1898bimetallism{]}. Bryan even invoked the
American Revolution explicitly: ``instead of having a gold standard
because England has, we shall restore bimetallism, and then let England
have bimetallism because the United States have''---a deliberate
re-enactment of 1776's monetary independence.

Bryan's peroration constitutes one of the most powerful expressions of
the critique of monometallism ever uttered in the American political
tradition:

\begin{quote}
``Having behind us the producing masses of this nation and the world,
supported by the commercial interests, the laboring interests, and the
toilers everywhere, we will answer their demand for a gold standard by
saying to them: `You shall not press down upon the brow of labor this
crown of thorns; you shall not crucify mankind upon a cross of
gold.'\,''
\end{quote}

The image of crucifixion upon gold directly inverts the Christian
symbolism that Blake deployed throughout his prophetic works: gold,
which should represent the solar principle of divine love, has been
weaponized into an instrument of torture---the miser's guinea elevated
above the sun until it becomes the cross upon which human labor is
sacrificed. Bryan's speech echoes Blake's demand in \emph{Milton}:
``Leave counting Gold \& diamonds, for Gods eye delights in thee''
(Plate 26, line 27; Erdman p.~121) {[}@erdman1988complete{]}. Yet Bryan
lost the 1896 election to William McKinley---the candidate of gold and
industrial capital---and the Gold Standard Act of 1900 formally
enthroned the monometallic regime, ending the American bimetallic
experiment that Hamilton had initiated 108 years earlier.
\end{frame}

\end{document}
